\section{Discussion and Limitations}

The central systems trade-off is separation. Building each native candidate on
one fixed typed sample and replaying its payload in a backend-local planner
overlay makes membership and order comparisons controlled: native payload
construction, ordinary statistics, workload, truth, and planner settings do
not change with every search move. This is why the mechanism is useful even
though the deployed estimator remains PostgreSQL's native estimator. The
separation is a design-time contract, not a measured end-to-end speedup over a
physical implementation.

The transfer boundary is equally important. Full-data deployment creates a new
native realization and can refresh ordinary statistics and externally managed
extended-statistics payloads. Across the four RQ2 datasets, sample-selected
designs retained utility in the controlled physical comparison, but paired
regressions remained and the magnitude varied. The sample-side objective does
not guarantee the full-data objective. Native \texttt{ANALYZE} Stability v2
offers only qualified post-hoc descriptive evidence: five realizations per
stratum showed no first-place winner reversal, while ordinary and extended
payload fingerprints varied. The v1 campaign remains failed/ineligible because
six dependency payloads were recorded empty without a direct SQL-NULL marker;
the reanalysis is not prospective confirmatory evidence and does not separate
ordinary-statistics from extended-statistics variation.

Recommendation order is part of the interface rather than an incidental SQL
rendering choice. The OID ablation shows that, under the fixed-payload
hypothetical contract, changing activation order can change Plan Rows for a
dataset-dependent subset of queries. The reference order is favorable among
the seven declared alternatives, but the result is not global ordering
optimality. Its synthetic overlapping-MCV witness is one bounded mechanism
diagnostic. Physical non-reference arms are operational observations because
their native payload and ordinary-statistics fingerprints also changed; they
cannot identify an OID-only effect.

The Advisor search claim is similarly conditional. Greedy ADD wins the
recorded fixed-$k$ comparison on Forest10, Power7, and DMV11, while Singleton
utility wins on Census13. Census13 and DMV11 Greedy traces contain
budget-censored feasible incumbents, and Forest10 uses a historical
3,600-second budget rather than the current 300-second Greedy/$K_s$ search
deadline. A configuration evaluation is not a uniform planner-call unit, so
the cancelled 2,000-evaluation experiment would not by itself have established
end-to-end resource fairness. These results support a transparent
planner-in-the-loop reference and a dataset-dependent benefit, not a globally
optimal algorithm or resource-normalized superiority.

For the completed RQ1b analysis, the evidence is narrower than a broad
generalization claim. It covers only four audited AreCEL datasets whose
validation and test workloads follow the same generator contract. Strict-unseen
removes exact source-query overlap, but not all structural similarity. The
reported outcome is planner cardinality-estimation q-error rather than query
execution latency, and paired regressions remain on every dataset. Improvement
magnitude is heterogeneous. Candidate IDs from independent advisor runs are
not established as physical-definition identifiers, so zero identifier overlap
does not prove disjoint PostgreSQL definitions. Census13 additionally retains
its documented historical baseline provenance and statistics-policy caveat.

Fourth, the remaining limitations are operational as well as scientific.
The four RQ2 datasets show a descriptive sample/full Spearman association of
0.478--0.968, but independent full-data realizations and distinct objectives
prevent a transfer guarantee. RQ3 exact Plan Rows agreement is limited to its
registered synthetic clause forms. RQ1b covers one audited workload generator;
strict-unseen removes exact source-query overlap, not all structural similarity,
and q-error is not query latency. External truth requires import and binding
validation, while production-exact truth can require expensive exact counts;
the absence of learned-model inference is therefore not zero operational cost.
Literature comparisons likewise are contextual rather than matched experiments.
Finally, stock deployment is not production readiness: existing statistics
remain outside Advisor ownership, a final \texttt{ANALYZE} can rebuild native
payloads, and the objective for $M^*$ is not a guarantee for
$E_{\mathrm{existing}}\cup M^*$. Reconciliation and refresh remain DBA and
operational concerns.
